%! Unit 2: Functions, Transformations, and Graph Analysis — the unit packet cover sheet.
%
% THIS FILE IS THE COVER. Both wrappers \input it — unit_cover/main.tex (student
% packet) and unit_cover_key/main.tex (key packet, which adds a page 2 of exam
% scoring notes) — so page 1 can never drift between the two packets. Edit the
% cover here, never in a wrapper.

% Full-bleed royal banner
\begin{tikzpicture}[remember picture, overlay]
  \fill[royal] (current page.north west)
    rectangle ([yshift=-0.9in]current page.north east);
\end{tikzpicture}
\vspace{-0.8in}

\begin{center}
  {\color{white}\textbf{\LARGE Algebra, Trigonometry, and Data Analysis}}\\[4pt]
  {\color{mistmid}\large\textbf{Unit 2: Functions, Transformations, and Graph Analysis}}\\[6pt]
  {\color{white}\textbf{\large Unit Overview \quad\textbar\quad Eight Lessons \quad\textbar\quad Unit Test}}
\end{center}

\vspace{0.20in}
\namedateperiod
\vspace{0.14in}

\noindent
\small
In Unit 1 you were handed an expression and asked to rewrite it. In Unit 2 you are handed a
\textbf{graph} and asked what it says. That is not a change of style --- it is what this course
asks for: in AFDA the characteristics of a function are investigated \emph{from the graph}. So
every lesson here is one more question you can ask a picture. What are its inputs and outputs?
What moved it off its parent? Where does it cross, turn, level off, or change rules? And, at the
end: given a table of numbers and no picture at all, \emph{which family does this even belong to?}

\vspace{0.12in}

% BOXGUARD: all three boxes below are breakable and all three fit on page 1 today
% with room to spare, so these guards are inert as written. They are sized to
% their own boxes so that if the cover ever grows, a box moves whole to a second
% page instead of splitting and stranding a stub — the guard is the cheap
% insurance, and each was verified not to cost this sheet its single page.
\boxguard[14]
\begin{tocbox}
\small
\renewcommand{\arraystretch}{1.32}
\begin{tabularx}{\linewidth}{@{} c >{\raggedright\arraybackslash}p{5.0cm} l X @{}}
\rowcolor{mist}
\textbf{\#} & \textbf{Lesson} & \textbf{SOL} & \textbf{The one idea} \\
\midrule
2.0 & Unit Opener                       & ---                       & A graph is read, not solved \\
2.1 & Notation, Domain, and Range       & \texttt{AF.2a, e}         & $f(3)=10$ and the point $(3,10)$ are one sentence \\
2.2 & Parent Functions \& Transformations & \texttt{AF.1a--b}       & Outside the parentheses, outputs; inside, inputs \\
2.3 & Equations and Graphs Both Ways    & \texttt{AF.1d, f}         & A matching picture is not proof \\
2.4 & Intercepts, Zeros, and Extrema    & \texttt{AF.2b--d}         & Say which axis the number came from \\
2.5 & End Behavior and Asymptotes       & \texttt{AF.2f--g}         & Getting closer and closer is not arriving \\
2.6 & Piecewise-Defined Functions       & \texttt{AF.2}             & Which piece owns this input? \\
2.7 & Choosing and Comparing Models     & \texttt{AF.1c, e, g; 2h}  & Chosen by how outputs \emph{change}, not by shape \\
\end{tabularx}
\end{tocbox}

\vspace{0.12in}

\boxguard[16]
\begin{spiralbox}
\small
\begin{enumerate}[leftmargin=1.5em, itemsep=3pt, topsep=2pt]
  \item \textbf{Before you write a number down, say which axis it came from.} An interval of
        \emph{inputs} and a stretch of \emph{outputs} can be the same two numbers, in the same
        units, and mean nothing alike. This is also why an extreme is reported as a
        \emph{value and a location} --- ``\$1800'' is half an answer.
  \item \textbf{Every transformation answers one question: did the outputs change, or did the
        inputs?} Outside the parentheses touches outputs; inside touches inputs. The horizontal
        rule is never ``the sign flips'' --- it is substitution: $f(x-3)$ at $x=5$ computes $f(2)$.
  \item \textbf{Getting closer and closer is not arriving.} That is why ``none'' is so often the
        right answer to ``what is the zero / the maximum?'' And an asymptote is a \emph{line}, so
        it gets an \emph{equation}: $y=20$, not $20$.
  \item \textbf{A rule owns only its own stretch of the domain.} A piecewise function is one
        function, not several; ask ``which piece owns this input?'' before anything else. A
        \emph{corner} is a change of rule, not a hole --- it is not a break.
  \item \textbf{A family is chosen by how the outputs \emph{change}, not by what the graph looks
        like.} ``It curves upward'' rules out linear and nothing else. Constant first differences
        mean linear, constant \emph{second} differences mean quadratic, a constant ratio means
        exponential --- over equal steps in $x$.
\end{enumerate}
\end{spiralbox}

\vspace{0.12in}

\boxguard[8]
\begin{remindbox}
\small
\textbf{The unit test has four parts and is worth 100 points:} Part A vocabulary matching
(8 pts), Part B multiple choice (12 pts), Part C short answer and computation (50 pts), and
Part D extended response (30 pts). A \textbf{practice test} with the same structure and
different numbers is bound at the back of this packet --- work it with no notes, then check it
against the key. Most of this test is a \emph{graph you have never seen}, so practice reading
one out loud: domain, range, intercepts, zeros, rising and falling, extremes with their
locations, and what happens at the far ends.
\end{remindbox}
